\section{Parton distributions}

\subsection{Generic matrix element  and parton distributions}

 The basic object for defining parton distributions is a matrix element
 of a bilocal operator that  (skipping inessential details of  its
 spin structure)   may be written   generically  like
 $\langle  p | \phi (0)  \phi(z)  | p \rangle $.
 Due to invariance under Lorentz transformations, it is given by a function of two scalars,  $(pz)$
 (which will  be denoted by $-\nu$)  and   $z^2$
 (or $-z^2$, in order to have a positive value  for \mbox{spacelike $z$)}
  \begin{align}
  \langle p |   \phi(0) \phi (z)|p \rangle %\equiv B(z,p)
=  & {\cal M} (-(pz), -z^2)  =  {\cal M} (\nu, -z^2)
\,  .
 \
 \label{lorentz}
\end{align}
 One can demonstrate \cite{Radyushkin:2016hsy,Radyushkin:1983wh}    that, for all relevant Feynman diagrams,
 its  Fourier transform  ${\cal P} (x, -z^2)$ with respect to $(pz)$
 has $-1 \leq x \leq 1$ as support, i.e.,
   \begin{align}
 {\cal M} (-(pz), -z^2)
&   =
 \int_{-1}^1 dx
 \, e^{-i x (pz) } \,  {\cal P} (x, -z^2)  \   .
  \label{MPD}
\end{align}
  Eq. (\ref{MPD}) serves as a covariant definition of $x$.
   In this  definition of $x$, one does not need to assume that
$p^2=0$ or $z^2=0$.

Choosing a light-like  $z$, e.g.,
having  solely the light-front component $z_-$,  we
parametrize the matrix element  by $f(x)$, the twist-2 parton distribution
  \begin{align}
  {\cal M} (-p_+ z_- , 0)  =
   \int_{-1}^1 dx \, f(x) \,
e^{-ixp_+ z_-} \,  \   .
 \label{twist2par0}
\end{align}
One can  rewrite this definition as
  \begin{align}
  {\cal M} (\nu, 0)  =
   \int_{-1}^1 dx \, f(x) \,
e^{ix\nu} \,  \  .
 \label{twist2par1}
\end{align}
The inverse relation is given by
\begin{align}
 f  (x)  =\frac{1}{2 \pi}  \int_{-\infty}^{\infty}  d\nu \, e^{-i x \nu}  \, {\cal M}(\nu , 0)    = {\cal P} (x, 0)
  \  .
\label{fxMnu}
\end{align}

Due to the fact that $f(x) =   {\cal P} (x, 0) $, the function $ {\cal P} (x, -z^2) $ provides a generalization of the concept of
PDFs onto non-lightlike intervals $z^2$ (in principle, $z^2$ may be even timelike).
 Following   \cite{Radyushkin:2017cyf} ,  we will be referring to it as the
 {\it pseudo-PDF}.  The variable $(pz)=-\nu $ is  called  often  the {\it Ioffe time} \cite{Ioffe:1969kf},
and  consequently ${\cal M}(\nu , -z^2) $  is  the {\it Ioffe-time distribution} \cite{Braun:1994jq}.


 In renormalizable theories (including QCD),
  the function  $ {\cal M} (\nu, -z^2)  $
 has  logarithmic $\sim \ln (-z^2)$  singularities
which generate the perturbative evolution of parton densities.
In the approach based on  the operator product expansion
(OPE),  the standard procedure is to  remove  these singularities  with the help of some
 prescription.
The most popular of them is the  $\overline{\rm MS}$  scheme based
on dimensional  regularization.
Consequently the resulting PDFs  have a dependence on the renormalization scale
$\mu$, and therefore one should write the PDFs as $   f(x, \mu^2)$.

At small spacelike $z^2$ and at the leading logarithm level, the pseudo-PDFs are
related to the $\overline{\rm MS}$ distributions by a simple rescaling of their second arguments.
In particular, when $z^2=-z_3^2$, one has
\begin{align}
    {\cal P} (x, z_3^2) = f \left (x, (2 e^{-\gamma_E}/ z_3)^2 \right )
  \  ,
\label{Ptof}
\end{align}
where $\gamma_E$  is the Euler's constant.
The rescaling  factor  between $\mu$ and $1/z_3$ is very close to 1, since \mbox{$2e^{-\gamma_E}=1.12$. }


\subsection{Transverse momentum dependent-  and quasidistributions}


Treating the target momentum $p$ as   longitudinal,
\mbox{$p= (E, {\bf 0}_\perp, P)$,}  one can introduce
transverse degrees of freedom.
In particular, taking $z$ that has
 $z_-$ and \mbox{$z_\perp= \{z_1,z_2\}$}  components only,  one defines  the
{\it TMD }   ${\cal F} (x, k_\perp^2)$   %as follows
      \begin{align}
 {\cal P} (x,  z_\perp^2)
&   =
    \int  d^2{\bf k}_\perp   e^{i ({\bf k}_\perp {\bf z}_\perp)}
      {\cal F} (x, {k}_\perp^2)
 \   .
  \label{MTMD0}
\end{align}
In  this context, the pseudo-PDFs  $ {\cal P} (x,  z_\perp^2) $ actually coincide
with the  {\it impact parameter distributions},   a familiar object
used in many TMD studies.


Since one cannot arrange  light-like separations on the lattice,
it was proposed  \cite{Ji:2013dva}  to consider equal-time
spacelike separations
$z= (0,0,0,z_3)$ (or, for brevity, \mbox{$z=z_3$}). Then,
 in  the   \mbox{$p=(E, 0_\perp, P)$}  frame,
 one  can introduce the quasi-PDF  $Q(y, P)$    through a parametrization
  \begin{align}
  \langle p |   \phi(0) \phi (z_3)|p \rangle %\equiv B(z,p)
=  &
\int_{-\infty}^{\infty}   dy \,
 Q(y, P) \,  e^{i y  P z_3 } \,
 \  .
 \label{newVDFxzQ}
\end{align}
According to  this definition, the  quasi-PDF  $Q(y,P)$ describes  the probability
that  the  parton carries the  fraction $y$ of the parent hadron's third momentum
component $P$.
The variables $\nu$ and $-z^2$
 in this case are given by $Pz_3$ and $z_3^2$, so we have
 \begin{align}
{\cal M} (\nu,z_3^2)
=  &
\int_{-\infty}^{\infty}   dy \,
 Q(y, P) \,  e^{i y  \nu } \,
 \  .
\end{align}
Since $z_3^2= \nu^2/P^2$,   the inverse  Fourier transformation
may be written as
\begin{align}
  Q(y,  P)   =\frac{1}{2 \pi}  \int_{-\infty}^{\infty}  d\nu \,
   \, e^{-i y  \nu}  \, {\cal M} (\nu,  \nu^2/P^2)    \  .
\label{IxM}
\end{align}
It
shows that $  Q(y,  P)  $  tends to $f(y)$  in the
$P \to \infty$  limit,  since  formally    ${\cal M} (\nu,  \nu^2/P^2) \to {\cal M} (\nu, 0)$
when \mbox{$P \to \infty$. }

As established  in Ref.  \cite{Radyushkin:2016hsy}, quasi-PDFs may be written in terms of TMDs %
 \begin{align}
 Q(y, P) /P =  & \,\int_{-\infty}^{\infty} d  k_1
\int_{-1} ^ {1} d  x \,   {\cal F} (x, k_1^2+(y-x)^2P^2 )
  \  .
 \label{QTMD}
 \end{align}

\subsection{Quantum chromodynamics (QCD)  case}

In case of %
  the non-singlet parton densities of QCD,
one is considering matrix elements
    \begin{align}
 {\cal M}^\alpha  (z,p) \equiv \langle  p |  \bar \psi (0) \,
 \gamma^\alpha \,  { \hat E} (0,z; A) \psi (z) | p \rangle \  ,
\label{Malpha}
\end{align}
 where  $
{ \hat E}(0,z; A)$ is  the  standard  $0\to z$ straight-line gauge link
 in the quark (fundamental) representation.
 These matrix elements can  be decomposed into $p^\alpha$ and $z^\alpha$ parts
\begin{align}
{\cal M}^\alpha  (z,p) = &2 p^\alpha  {\cal M}_p (-(zp), -z^2)
 + z^\alpha  {\cal M}_z (-(zp),-z^2)
\ .
\end{align}
Only %
the ${\cal M}_p (-(zp), -z^2) $ part gives the twist-2 distribution when  $z^2 \to 0$.

Introducing TMDs,   one  takes   $z=(z_-, z_\perp)$  and  the \mbox{$\alpha=+$}  component of
${\cal M}^\alpha$. Hence,  the  $z^\alpha$-part drops out.
After that,  $ {\cal M}_p(\nu, z_\perp^2)$ is   the only
surviving part of ${\cal M}^\alpha  (z,p)$, and in the remaining discussion
we use the short hand notation of ${\cal M} \equiv {\cal M}_p$.

In the case of    quasidistributions $Q(y,P)$,
we   can  avoid  the $z^\alpha$ contamination by   considering   the
  time component of \mbox{$ {\cal M}^\alpha  (z=z_3,p)$}  and defining
 \begin{align}
&  {\cal M}^0   (z_3,p)
  =  2p^0
 \int_{-1}^1 dy\,
Q(y,P) \,
 \,  e^{i  y Pz_3 }  \  .
 \label{OPhixspin12}
\end{align}



\subsection{Factorized models}

 The structure of the quasi-PDFs may be illustrated on the example
 of the simplest models in which
the nonperturbative (or soft) part of the TMDs
 ${\cal F}  (x, k_\perp^2)$ is  represented   by a product
    \begin{align}
{\cal F}^{\rm soft}  (x, k_\perp^2) = f(x) K(k_\perp^2)
\end{align}
 of the
 collinear parton distribution $f(x)$  and a \mbox{$k_\perp^2$-dependent}  factor $K(k_\perp^2)$, usually
 modeled by a Gaussian.
 As we shall see, the quasi-PDFs  have a rather complicated structure,
 even when they  are built from   these simple  factorized models.

 For  the Ioffe-time distribution  ${\cal M} (\nu,-z^2)$,  this Ansatz corresponds
to the factorization assumption
   \begin{align}
   {\cal M}^{\rm soft}  (\nu,z_3^2) = {\cal M}^{\rm soft}  (\nu,0){\cal M} (0,z_3^2) \
    \end{align}
    for its soft part.
Still, even if the soft TMD factorizes,  the soft part of the  quasi-PDF has the  convolution
structure of Eq. (\ref{QTMD}).    Taking,   for example,   a Gaussian  form
  \begin{align}
K_G(k_\perp^2) =  &
\frac{1}{\pi  \Lambda^2}  e^{-k_\perp^2/\Lambda^2}
  \  ,
 \label{DTMDsmallfac}
\end{align}
one gets the following model  for the quasi-PDF
 \begin{align}
 Q_G(y, P)  = &\frac{P}{\Lambda \sqrt{\pi} }  \,
 \int_{-1}^1 dx\,  %
f(x) \,
  e^{- (x -y)^2 P^2 / \Lambda^2 }
 \  .
 \label{QinG}
\end{align}
Choosing for $f(x)$
 a simple toy PDF resembling  the nucleon valence densities
$f(x)=4(1-x)^3 \theta (0\leq x \leq 1)$, one gets the curves shown in Fig.
 \ref{Qgy}.
 For large $P$, the quasi-PDF clearly tends to the $f(y)$   PDF  form.
However, only for   %%
$P \sim 10 \Lambda$  one gets  a quasi-PDF that is rather close to the $P \to \infty$ limiting shape.
Still, since  $\Lambda \sim \langle k_\perp \rangle$,
one   translates  the
 \mbox{$P\sim 10 \Lambda$}  estimate into $P \sim 3$ GeV,
which is  rather  large.


 \begin{figure}[t]
    \centerline{\includegraphics[width=3.6in]{images/QyPn.pdf}}
    \caption{Evolution of quasi-PDF $Q(y,P)$  in the  factorized Gaussian  model  for {$P/\Lambda =1, 5, 10, 50$}.
    \label{Qgy}}
    \end{figure}
